\documentclass[a4paper,12pt]{article}
\usepackage[utf8]{inputenc}
\usepackage[T1]{fontenc}
\usepackage[french]{babel}
\usepackage{geometry}
\usepackage{graphicx}
\usepackage{fancyhdr}
\usepackage{lastpage}
\usepackage{xcolor}
\usepackage{enumitem}
\usepackage{hyperref}
\usepackage{titlesec}
\usepackage{parskip}

% Page layout amélioré avec des marges plus grandes
\geometry{
    margin=3cm,
    top=4cm,       % Marge supérieure plus grande pour l'en-tête
    bottom=3.5cm,  % Marge inférieure plus grande pour le pied de page
    headheight=2cm % Hauteur d'en-tête adaptée au logo
}

\pagestyle{fancy}
\fancyhf{}
\fancyhead[L]{\includegraphics[height=1.5cm]{smblogo.png}} % Hauteur fixe pour contrôler l'espace
\fancyhead[R]{\textbf{SMB Le Bon Point}}
\fancyfoot[L]{\footnotesize SMB Le Bon Point -- SIREN 908247836 -- contact@smb-lebonpoint.fr}
\fancyfoot[R]{\footnotesize Page \thepage\ sur \pageref{LastPage}}

% Espacements supplémentaires pour les en-têtes et pieds de page
\setlength{\footskip}{30pt}

% Couleurs
\definecolor{primary}{RGB}{33, 37, 41}
\definecolor{secondary}{RGB}{99, 102, 241}

% Format des sections avec espacement amélioré
\titleformat{\section}{\large\bfseries\color{primary}}{\thesection}{1em}{}
\titleformat{\subsection}{\normalsize\bfseries\color{primary}}{\thesubsection}{1em}{}

% Ajout d'espacement entre les sections et sous-sections
\titlespacing*{\section}{0pt}{3.5ex plus 1ex minus .2ex}{2.3ex plus .2ex}
\titlespacing*{\subsection}{0pt}{3.25ex plus 1ex minus .2ex}{1.5ex plus .2ex}

% Configuration des listes avec plus d'espacement
\setlist{noitemsep, topsep=6pt, parsep=0pt, partopsep=0pt}


% Titre sans date
\title{\textbf{Conditions Générales de Vente -- SMB Le Bon Point}}
\author{}
\date{}  % Date vide pour ne pas l'afficher

\begin{document}

\maketitle
\vspace{-1.5cm} % Réduit légèrement mais pas trop

\begin{center}


    SARL au capital de [montant en euros] \\
    \vspace{0.2cm} % Ajout d'espacement
    
    Siège social : [Adresse complète, ex. 12 Rue de Paris, 94100 Saint-Maur-des-Fossés] \\
    SIREN : 908247836 -- RCS Créteil \\
    Téléphone : 01 23 45 67 89 \\
    Email : \href{mailto:contact@smb-lebonpoint.fr}{contact@smb-lebonpoint.fr} \\
    Site web : \href{https://www.smb-lebonpoint.fr}{www.smb-lebonpoint.fr}
\end{center}
\vspace{0.5cm} % Espace avant le préambule

\section*{Préambule}
Les présentes Conditions Générales de Vente (CGV) régissent l'ensemble des relations contractuelles entre la société SMB Le Bon Point, ci-après dénommée « la Société », et toute personne physique ou morale, ci-après dénommée « le Stagiaire », souhaitant s'inscrire à un stage de sensibilisation à la sécurité routière organisé par la Société. Ces stages, conformes aux dispositions du Code de la route et sous agrément préfectoral, permettent de récupérer des points sur le permis de conduire ou de répondre à une obligation légale, telle que définie par l'article L223-6 du Code de la route. 

Les CGV s'appliquent sans restriction ni réserve à toutes les inscriptions effectuées via le site internet \href{https://www.smb-lebonpoint.fr}{www.smb-lebonpoint.fr}, par téléphone, ou tout autre canal autorisé par la Société. En s'inscrivant, le Stagiaire reconnaît avoir pris connaissance et accepté les présentes CGV dans leur intégralité, ainsi que les éventuelles conditions particulières mentionnées sur le site.

\section{Définitions}
Pour l'interprétation des présentes CGV, les termes suivants sont définis comme suit :
\begin{itemize}
    \item \textbf{Société} : SMB Le Bon Point, entité organisatrice des stages, identifiée par le SIREN 908247836.
    \item \textbf{Stagiaire} : Toute personne physique ou morale s'inscrivant à un stage.
    \item \textbf{Stage} : Formation de sensibilisation à la sécurité routière, agréée par la préfecture, visant la récupération de points ou le respect d'une obligation légale.
    \item \textbf{Site} : Plateforme en ligne \href{https://www.smb-lebonpoint.fr}{www.smb-lebonpoint.fr}, utilisée pour les inscriptions et paiements.
    \item \textbf{ANTS} : Agence Nationale des Titres Sécurisés, autorité délivrant les numéros d'agrément des stages.
\end{itemize}

\section{Objet}
Les présentes CGV ont pour objet de définir les droits et obligations des parties dans le cadre de l'inscription, de l'organisation, et de la participation aux stages de sensibilisation à la sécurité routière proposés par la Société. Elles précisent les conditions de vente, de paiement, d'annulation, et de traitement des données personnelles, ainsi que les responsabilités respectives de la Société et du Stagiaire. 

Les stages sont dispensés dans le respect des normes fixées par le Code de la route, sous la supervision d'un formateur et d'un psychologue agréés, et visent à promouvoir une conduite responsable et sécurisée.

\section{Inscription}
\subsection{Modalités d'inscription}
L'inscription à un stage s'effectue principalement via le formulaire en ligne disponible sur \href{https://www.smb-lebonpoint.fr}{www.smb-lebonpoint.fr}. Le Stagiaire doit :
\begin{itemize}
    \item Sélectionner un stage parmi les sessions proposées, indiquant le lieu, les dates, et le prix.
    \item Remplir un formulaire électronique avec les informations suivantes : civilité, nom, prénom, adresse complète, code postal, ville, numéro de téléphone, adresse email, nationalité, date et lieu de naissance, numéro de permis de conduire, date de délivrance du permis, préfecture de délivrance, état du permis (valide ou invalide), et motif du stage (ex. : stage volontaire ou obligatoire).
    \item Téléverser, le cas échéant, les documents suivants au format numérique via la plateforme sécurisée Supabase : recto et verso du permis de conduire (\texttt{permis\_recto}, \texttt{permis\_verso}), recto et verso d'une pièce d'identité (\texttt{id\_recto}, \texttt{id\_verso}), lettre 48N (\texttt{letter\_48N}), et tout document supplémentaire (\texttt{extraDocument}).
    \item Confirmer l'exactitude des informations fournies et accepter les présentes CGV en cochant la case prévue à cet effet.
\end{itemize}

Le Stagiaire est tenu de vérifier l'exactitude des informations saisies, car toute erreur peut entraîner un refus de validation du stage par la préfecture compétente. La Société ne saurait être tenue responsable des conséquences découlant d'informations erronées ou incomplètes.

\subsection{Validation de l'inscription}
L'inscription est considérée comme définitive après :
\begin{itemize}
    \item Validation complète du formulaire en ligne par le Stagiaire.
    \item Paiement intégral du prix du stage via la plateforme sécurisée Stripe, intégrée au Site.
    \item Envoi d'un email de confirmation au Stagiaire, contenant les informations suivantes : numéro de stage ANTS, dates et lieu du stage, horaires, et une attestation d'inscription au format PDF stockée dans la base de données de la Société.
\end{itemize}

La Société se réserve le droit de refuser une inscription en cas de non-respect des conditions d'éligibilité (ex. : permis de conduire invalide sans justificatif) ou de non-paiement.

\subsection{Capacité des stages}
Chaque stage est limité à un maximum de 20 participants, sauf indication contraire sur le Site. Les inscriptions sont acceptées dans l'ordre d'arrivée, dans la limite des places disponibles. En cas de dépassement de la capacité, la Société peut proposer une inscription à une autre session ou refuser l'inscription sans engager sa responsabilité.

\section{Prix et Modalités de Paiement}
\subsection{Prix}
Le prix de chaque stage, exprimé en euros toutes taxes comprises (TTC), est clairement indiqué sur \href{https://www.smb-lebonpoint.fr}{www.smb-lebonpoint.fr} lors de la sélection du stage. Ce prix varie en fonction du lieu, des dates, et de la nature du stage (volontaire, obligatoire en période probatoire, ou autre). 

Aucun frais supplémentaire ne sera facturé au Stagiaire, sauf en cas de demande expresse de services additionnels (ex. : envoi postal de documents). Les prix sont susceptibles d'évoluer, mais le prix applicable est celui affiché au moment de l'inscription.

\subsection{Paiement}
Le paiement est effectué exclusivement en ligne via la plateforme sécurisée Stripe, intégrée au Site. Le Stagiaire doit régler l'intégralité du prix du stage au moment de l'inscription. Les moyens de paiement acceptés incluent les cartes bancaires (Visa, Mastercard, American Express, etc.). 

La Société ne conserve ni n'accède aux données bancaires du Stagiaire, celles-ci étant traitées directement par Stripe conformément aux normes de sécurité PCI-DSS. En cas d'échec du paiement, l'inscription est annulée, et le Stagiaire est informé par email.

\subsection{Facturation}
Une facture électronique est émise et envoyée au Stagiaire par email après confirmation du paiement. Cette facture détaille le numéro de stage ANTS, le prix payé, les dates et le lieu du stage, ainsi que les coordonnées de la Société. Sur demande expresse, une facture papier peut être envoyée par courrier, sous réserve de frais supplémentaires.

\section{Droit de Rétractation}
Conformément à l'article L221-18 du Code de la consommation, le Stagiaire dispose d'un délai de 14 jours calendaires à compter de la réception de l'email de confirmation pour exercer son droit de rétractation, sans avoir à justifier de motif ni à supporter de pénalités, sauf si le stage a débuté avant l'expiration de ce délai. 

En application de l'article L221-28 du Code de la consommation, le droit de rétractation ne s'applique pas si le Stagiaire a expressément consenti à l'exécution immédiate du service (ex. : participation à un stage commençant dans moins de 14 jours) et a renoncé à ce droit lors de l'inscription.

Pour exercer son droit de rétractation, le Stagiaire doit adresser une notification claire et non équivoque à la Société, par email à \href{mailto:contact@smb-lebonpoint.fr}{contact@smb-lebonpoint.fr} ou par courrier recommandé avec accusé de réception à [adresse complète], en utilisant, si souhaité, le modèle de formulaire de rétractation fourni en annexe. La Société remboursera le Stagiaire de la totalité des sommes versées, y compris les frais de livraison éventuels, dans un délai de 14 jours à compter de la réception de la notification, par le même moyen de paiement utilisé initialement.

\section{Annulation et Remboursement}
\subsection{Annulation par le Stagiaire}
Le Stagiaire peut annuler son inscription selon les conditions suivantes :
\begin{itemize}
    \item \textbf{Dans le délai de rétractation} : Si le droit de rétractation s'applique, le Stagiaire est remboursé intégralement, sans frais, sous 14 jours.
    \item \textbf{Hors délai de rétractation} :
    \begin{itemize}
        \item Annulation notifiée plus de 15 jours calendaires avant le début du stage : remboursement de 50 \% du prix du stage, déduction faite des frais administratifs.
        \item Annulation notifiée moins de 15 jours avant le début du stage : aucun remboursement, sauf en cas de force majeure dûment justifiée (ex. : certificat médical, décès d'un proche), auquel cas un remboursement partiel ou un report peut être envisagé à la discrétion de la Société.
    \end{itemize}
\end{itemize}

Toute demande d'annulation doit être adressée par email à \href{mailto:contact@smb-lebonpoint.fr}{contact@smb-lebonpoint.fr}, avec les informations suivantes : nom, prénom, numéro de stage ANTS, et motif de l'annulation. La Société accuse réception de la demande dans un délai de 48 heures.

\subsection{Annulation par la Société}
La Société se réserve le droit d'annuler un stage dans les cas suivants :
\begin{itemize}
    \item \textbf{Force majeure} : Événements imprévisibles et indépendants de la volonté de la Société (ex. : catastrophe naturelle, restrictions gouvernementales).
    \item \textbf{Insuffisance de participants} : Nombre de Stagiaires inscrits inférieur au minimum requis (6 participants).
    \item \textbf{Indisponibilité} : Absence imprévue de l'instructeur ou du psychologue agréé.
\end{itemize}

En cas d'annulation, la Société informe le Stagiaire par email ou téléphone dans les plus brefs délais, au minimum 48 heures avant le début prévu du stage, sauf en cas de force majeure. Le Stagiaire peut choisir :
\begin{itemize}
    \item Un report sans frais sur une autre session de son choix, sous réserve de disponibilité.
    \item Un remboursement intégral du prix payé, effectué sous 14 jours par le même moyen de paiement.
\end{itemize}

La Société ne saurait être tenue responsable des frais annexes engagés par le Stagiaire (ex. : transport, hébergement) en cas d'annulation.

\section{Obligations des Parties}
\subsection{Obligations du Stagiaire}
Le Stagiaire s'engage à :
\begin{itemize}
    \item Fournir des informations complètes, exactes, et à jour lors de l'inscription, notamment en ce qui concerne son identité, son permis de conduire, et son éligibilité au stage.
    \item Téléverser des documents lisibles et valides via la plateforme Supabase, au format JPEG, PNG, ou PDF, conformément aux exigences du formulaire (ex. : \texttt{permis\_recto}, \texttt{id\_recto}, \texttt{letter\_48N}).
    \item Respecter les horaires du stage, qui se déroule sur deux jours consécutifs, avec une présence obligatoire pour l'ensemble des sessions (durée totale : environ 14 heures).
    \item Adopter un comportement respectueux envers les formateurs, les psychologues, et les autres participants, sous peine d'exclusion sans remboursement.
    \item Présenter une pièce d'identité en cours de validité et, le cas échéant, son permis de conduire ou la lettre 48N, le premier jour du stage.
\end{itemize}

Tout manquement à ces obligations peut entraîner le refus d'accès au stage ou l'invalidation de l'attestation de suivi, sans possibilité de remboursement.

\subsection{Obligations de la Société}
La Société s'engage à :
\begin{itemize}
    \item Organiser les stages dans le respect des dispositions réglementaires, avec des formateurs et psychologues disposant d'un numéro d'autorisation préfectorale.
    \item Mettre à disposition un lieu de formation adapté, conforme aux exigences de sécurité et d'accessibilité.
    \item Délivrer une attestation de suivi de stage au Stagiaire ayant participé intégralement, et transmettre cette attestation à la préfecture compétente pour la récupération de points ou la validation de l'obligation légale.
    \item Assurer la confidentialité et la sécurité des données personnelles et documents téléversés, conformément au RGPD.
    \item Informer le Stagiaire de tout changement relatif au stage (ex. : modification d'horaire, annulation) dans les meilleurs délais.
\end{itemize}

\section{Protection des Données Personnelles}
Conformément au Règlement Général sur la Protection des Données (RGPD) et à la loi Informatique et Libertés, la Société collecte et traite les données personnelles du Stagiaire dans le but de :
\begin{itemize}
    \item Gérer les inscriptions, les paiements, et l'organisation des stages.
    \item Émettre et transmettre les attestations de suivi aux autorités préfectorales.
    \item Communiquer avec le Stagiaire (confirmation d'inscription, facture, informations sur le stage).
\end{itemize}

Les données collectées incluent : nom, prénom, adresse, téléphone, email, nationalité, date et lieu de naissance, numéro de permis, et documents téléversés (\texttt{permis\_recto}, \texttt{id\_recto}, etc.). Les documents sont stockés sur la plateforme sécurisée Supabase, et les autres données dans une base PostgreSQL conforme aux normes de sécurité. Ces données ne sont partagées qu'avec les autorités compétentes (ex. : préfecture) lorsque requis par la loi.

Le Stagiaire dispose des droits suivants sur ses données :
\begin{itemize}
    \item Droit d'accès : Consulter les données détenues.
    \item Droit de rectification : Corriger les données inexactes.
    \item Droit de suppression : Demander la suppression, sous réserve des obligations légales de conservation.
    \item Droit à la portabilité : Recevoir les données dans un format structuré.
    \item Droit d'opposition : S'opposer au traitement pour des raisons légitimes.
\end{itemize}

Ces droits peuvent être exercés en adressant une demande écrite à \href{mailto:contact@smb-lebonpoint.fr}{contact@smb-lebonpoint.fr}, accompagnée d'une copie d'une pièce d'identité. La Société répondra dans un délai d'un mois, sauf cas exceptionnel. En cas de réclamation, le Stagiaire peut saisir la CNIL (Commission Nationale de l'Informatique et des Libertés, \href{https://www.cnil.fr}{www.cnil.fr}).

\section{Responsabilité}
La Société ne peut être tenue responsable dans les cas suivants :
\begin{itemize}
    \item Inexactitude ou omission des informations fournies par le Stagiaire, entraînant un refus de validation du stage par la préfecture ou une impossibilité de participer.
    \item Non-présentation du Stagiaire au stage, retard, ou non-respect des conditions de participation (ex. : absence de documents requis).
    \item Problèmes techniques indépendants de la volonté de la Société, tels que des interruptions temporaires du Site, de la plateforme Stripe, ou du service Supabase, bien que toutes les mesures raisonnables soient prises pour garantir la continuité du service.
    \item Préjudices indirects subis par le Stagiaire, tels que des frais de déplacement ou d'hébergement, en cas d'annulation ou de report du stage.
\end{itemize}

La responsabilité de la Société est strictement limitée au prix du stage payé par le Stagiaire, sauf en cas de faute lourde ou intentionnelle.

\section{Force Majeure}
Conformément à l'article 1218 du Code civil, aucune des parties ne peut être tenue responsable d'un manquement à ses obligations en cas de force majeure, définie comme un événement imprévisible, irrésistible, et extérieur (ex. : catastrophe naturelle, grève générale, pandémie). 

En cas de force majeure empêchant la tenue d'un stage, la Société informera le Stagiaire et proposera un report sur une autre session ou un remboursement intégral, conformément à l'article 6.2. Les frais annexes (transport, hébergement) ne seront pas remboursés.

\section{Assurance}
La Société souscrit une assurance responsabilité civile professionnelle couvrant les dommages causés aux Stagiaires ou à des tiers dans le cadre de l'organisation des stages. Les Stagiaires sont tenus de vérifier que leur propre assurance (ex. : responsabilité civile) couvre leur participation au stage. 

La Société ne peut être tenue responsable des dommages matériels ou corporels causés par un Stagiaire à un autre participant ou à un tiers pendant le stage.

\section{Litiges et Médiation}
En cas de différend relatif à l'interprétation ou à l'exécution des présentes CGV, le Stagiaire est invité à contacter la Société par email (\href{mailto:contact@smb-lebonpoint.fr}{contact@smb-lebonpoint.fr}) ou par courrier à [adresse complète] pour rechercher une solution amiable. La Société s'engage à répondre dans un délai de 15 jours ouvrables.

Conformément à l'article L616-1 du Code de la consommation, le Stagiaire peut recourir à une médiation gratuite en cas d'échec de la résolution amiable. Les coordonnées du médiateur de la consommation agréé seront communiquées sur demande ou disponibles sur \href{https://www.smb-lebonpoint.fr}{www.smb-lebonpoint.fr}. La médiation doit être saisie dans un délai d'un an à compter de la réclamation initiale.

À défaut de résolution amiable ou par médiation, tout litige sera soumis aux tribunaux compétents du ressort de Créteil, sauf disposition légale contraire imposant une autre juridiction. Les présentes CGV sont régies par le droit français.

\section{Dispositions Diverses}
\subsection{Modification des CGV}
La Société se réserve le droit de modifier les présentes CGV à tout moment pour refléter des évolutions légales, réglementaires, ou commerciales. Les CGV applicables sont celles en vigueur au moment de l'inscription du Stagiaire, disponibles sur \href{https://www.smb-lebonpoint.fr}{www.smb-lebonpoint.fr}. Le Stagiaire est invité à consulter régulièrement le Site pour prendre connaissance des éventuelles modifications.

\subsection{Nullité partielle}
Si une disposition des présentes CGV est déclarée nulle ou inapplicable par une autorité compétente, les autres dispositions resteront en vigueur dans la mesure du possible. Les parties s'efforceront de remplacer la disposition nulle par une clause valide ayant un effet économique similaire.

\subsection{Langue du contrat}
Les présentes CGV sont rédigées en français. Toute traduction fournie à titre informatif n'a pas de valeur juridique. En cas de divergence entre la version française et une traduction, la version française prévaut.

\subsection{Non-renonciation}
Le fait pour l'une des parties de ne pas exiger l'application d'une disposition des présentes CGV ne constitue pas une renonciation à ses droits ni une modification du contrat.

\vspace{15cm} % Espace avant l'annexe

\section*{Annexe : Modèle de Formulaire de Rétractation}
\begin{quote}
    \textbf{À l'attention de :} \\
    SMB Le Bon Point \\
    [Adresse complète, ex. 12 Rue de Paris, 94100 Saint-Maur-des-Fossés] \\
    Email : \href{mailto:contact@smb-lebonpoint.fr}{contact@smb-lebonpoint.fr} \\
    
    \textbf{Objet : Rétractation du contrat d'inscription à un stage} \\
    
    Je soussigné(e), [Nom et prénom du Stagiaire], notifie par la présente ma rétractation du contrat relatif à l'inscription au stage de sensibilisation à la sécurité routière portant le numéro ANTS [numéro du stage], prévu aux dates [dates du stage]. \\
    
    \textbf{Date d'inscription} : [date de l'inscription] \\
    \textbf{Nom et prénom} : [nom et prénom] \\
    \textbf{Adresse complète} : [adresse du Stagiaire] \\
    \textbf{Email} : [email du Stagiaire] \\
    
    Conformément à l'article L221-18 du Code de la consommation, je demande le remboursement intégral des sommes versées dans un délai de 14 jours. \\
    
    \textbf{Date} : [date de la demande] \\
    \textbf{Signature} : [uniquement pour une demande papier]
\end{quote}

\end{document}